\documentclass[a4paper,11pt]{article}
\usepackage[T1]{fontenc}
\usepackage[british]{babel}
\usepackage[margin=24mm]{geometry}
\usepackage{graphicx}
\title{Cadenza: local evaluation evidence}
\author{}
\date{14 October 2026}
\begin{document}
\maketitle
% Generated from repaired-version retained raw data by scripts/paper.py.
\section{Experimental evaluation}
\label{sec:evaluation}
We evaluated the repaired articulation monitor, its preserved assessment
boundaries and local execution costs. Service-path performances used
deterministic synthetic MIDI; direct checks used synthetic enriched traces.
No pianist recordings or user study were involved.
The previous evaluation, expectations and diagnostics were archived before
repair. Content hashes identify both source versions because usable project
Git history was unavailable; no parent repository commit is attributed to Cadenza.

\paragraph{Repair and behavioural evidence.}
Replaying the original final-marked-note input reproduced
\texttt{type\_error(clpr\_expression,null)}. The final attack has no next
distinct performed onset, so enrichment intentionally supplies a null gap.
The RML tried a numeric duration alternative before the existing no-gap
alternative, exposing null to generated CLP(R) arithmetic. We moved the
no-gap alternative first and regenerated Prolog with the vendored compiler.
Staccato still requires $d\leq0.75g$, and tenuto/legato require
$d\geq0.75g$ whenever the gap applies. No successor interval was invented.
Acceptance without a gap means that this requirement is inapplicable,
not that the final duration was verified.

The unchanged 28-case suite uses an original eight-note MusicXML
exercise with a chord, dynamics, staccato, tenuto and a pedal span. A separately
versioned manual ledger was frozen before repaired execution; it changes only
the expected final-note articulation error to no-gap acceptance.
All 196 metric outcomes were compared with the archive:
only this articulation verdict changed (exit 2 to 0); the other
195 were unchanged. All repaired predictions matched,
with 0 service failures and 0 monitor errors.
A further 29 focused checks also matched expectations:
18 direct compiled-monitor fixtures and 11 complete
conversion/service fixtures. They cover marked/unmarked null-gap finals,
durations below/at/above the 0.75 boundary, and earlier violations followed
by no-gap events; earlier violations still reject. The importer produces
staccato and tenuto; literal legato is tested directly at the monitor boundary.
Two malformed numeric fixtures retain their expected monitor errors;
the repair does not convert arbitrary arithmetic failures into acceptance.

\paragraph{Remaining assessment limits.}
Missing and duplicate-note cases still pass all seven metrics, showing the
absence of general score-coverage and one-to-one matching obligations.
Pedal endpoint association remains narrower than the nominal RML assessment
tolerance: this 120 BPM fixture accepts +25\,ms and rejects +26\,ms,
despite a nominal 0.25-beat (125\,ms) allowance. The effective boundary depends
on enrichment, tempo and floating-point snapping. Syncopated-pedal checks
event order without a numerical deadline: a 300\,ms delayed press passes
that metric, although pedal-toggle rejects it. These policies were unchanged.

\paragraph{Local execution costs.}
The same repository-score inputs yielded 216 and
2,324 parsed notes and 464 and
4,946 enriched events; hashes and counts match the archive.
The larger score's 2,418 pitched XML
noteheads include 94 tied continuations.
Onsets depend on the implementation's own timing map; this does not
independently validate alignment. Synthetic timestamps already have the
intended score-time origin, leaving live capture/playback clock alignment untested.

Each workload received three warm-ups and 30 measured calls, sequentially
in one Python process, with a new Prolog process per request and all seven
checks. A monotonic clock surrounded the actual \texttt{analyse} call,
including validation, temporary score creation, MusicXML parsing, enrichment,
JSONL serialisation/writing, Prolog startup/specification loading, monitoring,
verdict decoding and cleanup. Browser capture/conversion, MIDI decoding,
HTTP transport, rendering, Python startup/import and ahead-of-time compilation
were excluded. Stage-timer and warning-capture overhead remained included;
totals were measured directly, rather than summing component medians.

\begin{table}[t]
\centering
\resizebox{\linewidth}{!}{% Generated by scripts/summarise.py; all durations are seconds.
\begin{tabular}{lrrrrrrrrr}
\hline
Workload & Notes & Events & Warm & Runs & Median & Q1 & Q3 & Min & Max \\
\hline
score-2324 & 2324 & 4946 & 3 & 30 & 6.037239 & 5.956440 & 6.375752 & 5.754788 & 6.750570 \\
score-216 & 216 & 464 & 3 & 30 & 0.779284 & 0.758785 & 0.799557 & 0.726207 & 0.868918 \\
\hline
\end{tabular}
}
\caption{New local analysis times in seconds. Quartiles use linear
interpolation; all measured samples are retained and warm-ups excluded.}
\label{tab:local-cost}
\end{table}

Median totals were 0.779284\,s and 6.037239\,s
(Table~\ref{tab:local-cost}). Across 60 measured calls,
there were 0 service failures, 0 timeouts and
0 calls with monitor errors. The other six metrics accepted
on every call, but syncopated-pedal rejected on both workloads.
Separate standalone-monitor diagnostics locate rejection at events
28/464 and
29/4946,
respectively: at adjacent pedal spans the preceding region consumes a
boundary attack, leaving the next press without a fresh preceding attack.
These diagnostics were excluded from timing. Rejecting monitors can stop
early, so these costs do not represent seven fully accepting checks or
worst-case trace consumption. Parsing warnings and Prolog diagnostics were
retained with every sample; earlier timings were not reused as samples or targets.

Measurements ran on 14 October 2026 using
Intel(R) Core(TM) i7-7700HQ CPU @ 2.80GHz, 15.51\,GiB reported RAM, Linux 6.8.0-138-generic,
Python 3.11.0rc1, Partitura 1.9.0, NumPy 1.26.4, SciPy 1.17.1, lxml 6.1.3, mido 1.3.3 and
SWI-Prolog version 9.1.21 for x86\_64-linux. Numerical-library thread limits were
one; Python libraries and OS caches could remain warm. CPU frequency and
background load were uncontrolled. Exact versions, commands, hashes, raw
inputs/traces, manual expectations, all verdicts and timing samples accompany
the report. The evidence supports the repaired no-gap policy and these local
costs; it establishes neither complete musical assessment, independent live
alignment accuracy, scalability, real-time bounds nor user-facing latency.

\clearpage
\section*{Detailed behavioural results}
Cells are expected/actual exit codes: 0 acceptance, 1 rejection, 2 error.
A: articulation; C: chord jitter; D: dynamics; T: pedal toggle;
P: pitch agreement; R: pitch range; S: syncopated pedal.
Teacher-error labels follow the declared synthetic exercise rubric.
\begin{center}
\resizebox{\textwidth}{!}{% Generated by scripts/summarise.py. Cells are expected/actual exit codes.
% A articulation; C chord jitter; D dynamics; T pedal toggle;
% P pitch agreement; R pitch range; S syncopated pedal.
\begin{tabular}{llrrrrrrrl}
\hline
Case & Teacher error & A & C & D & T & P & R & S & Diff. \\
\hline
baseline & no & 0/0 & 0/0 & 0/0 & 0/0 & 0/0 & 0/0 & 0/0 & no \\
wrong\_pitch & yes & 0/0 & 0/0 & 0/0 & 0/0 & 1/1 & 0/0 & 0/0 & no \\
outside\_range & yes & 0/0 & 0/0 & 0/0 & 0/0 & 1/1 & 1/1 & 0/0 & no \\
bad\_velocity & yes & 0/0 & 0/0 & 1/1 & 0/0 & 0/0 & 0/0 & 0/0 & no \\
bad\_articulation & yes & 1/1 & 0/0 & 0/0 & 0/0 & 0/0 & 0/0 & 0/0 & no \\
chord\_spread\_29ms & no & 0/0 & 0/0 & 0/0 & 0/0 & 0/0 & 0/0 & 0/0 & no \\
chord\_spread\_30ms & yes & 0/0 & 1/1 & 0/0 & 0/0 & 0/0 & 0/0 & 0/0 & no \\
chord\_spread\_31ms & yes & 0/0 & 1/1 & 0/0 & 0/0 & 0/0 & 0/0 & 0/0 & no \\
pedal\_press\_24ms & no & 0/0 & 0/0 & 0/0 & 0/0 & 0/0 & 0/0 & 0/0 & no \\
pedal\_press\_25ms & no & 0/0 & 0/0 & 0/0 & 0/0 & 0/0 & 0/0 & 0/0 & no \\
pedal\_press\_26ms & no & 0/0 & 0/0 & 0/0 & 1/1 & 0/0 & 0/0 & 0/0 & no \\
pedal\_press\_124ms & no & 0/0 & 0/0 & 0/0 & 1/1 & 0/0 & 0/0 & 0/0 & no \\
pedal\_press\_125ms & no & 0/0 & 0/0 & 0/0 & 1/1 & 0/0 & 0/0 & 0/0 & no \\
pedal\_press\_126ms & yes & 0/0 & 0/0 & 0/0 & 1/1 & 0/0 & 0/0 & 0/0 & no \\
pedal\_release\_late\_25ms & no & 0/0 & 0/0 & 0/0 & 0/0 & 0/0 & 0/0 & 0/0 & no \\
pedal\_release\_late\_26ms & no & 0/0 & 0/0 & 0/0 & 1/1 & 0/0 & 0/0 & 0/0 & no \\
missing\_note & yes & 0/0 & 0/0 & 0/0 & 0/0 & 0/0 & 0/0 & 0/0 & no \\
duplicate\_note & yes & 0/0 & 0/0 & 0/0 & 0/0 & 0/0 & 0/0 & 0/0 & no \\
missing\_release & yes & 0/0 & 0/0 & 0/0 & 1/1 & 0/0 & 0/0 & 0/0 & no \\
pre\_attack & yes & 0/0 & 0/0 & 0/0 & 0/0 & 0/0 & 0/0 & 1/1 & no \\
pedal\_all\_absent & yes & 0/0 & 0/0 & 0/0 & 0/0 & 0/0 & 0/0 & 1/1 & no \\
late\_press\_300ms & yes & 0/0 & 0/0 & 0/0 & 1/1 & 0/0 & 0/0 & 0/0 & no \\
staccato\_375ms & no & 0/0 & 0/0 & 0/0 & 0/0 & 0/0 & 0/0 & 0/0 & no \\
tenuto\_374ms & yes & 1/1 & 0/0 & 0/0 & 0/0 & 0/0 & 0/0 & 0/0 & no \\
tenuto\_375ms & no & 0/0 & 0/0 & 0/0 & 0/0 & 0/0 & 0/0 & 0/0 & no \\
mf\_velocity\_64 & no & 0/0 & 0/0 & 0/0 & 0/0 & 0/0 & 0/0 & 0/0 & no \\
mf\_velocity\_80 & yes & 0/0 & 0/0 & 1/1 & 0/0 & 0/0 & 0/0 & 0/0 & no \\
final\_marked\_note & no & 0/0 & 0/0 & 0/0 & 0/0 & 0/0 & 0/0 & 0/0 & no \\
\hline
\end{tabular}
}
\end{center}
\clearpage
\section*{Focused articulation regressions}
These additional fixtures remain separate from the 28-case suite.
Cells show expected/actual articulation exit codes for the archived original
and repaired monitor. Direct fixtures run articulation only; complete service
fixtures run all seven checks, with the other six accepting in both versions.
No-gap acceptance means the duration requirement is inapplicable.
\begin{center}
\resizebox{\textwidth}{!}{% Generated by scripts/regression_report.py directly from results/regressions.json.
% Cells are expected/actual articulation exit codes: 0 accept, 1 reject, 2 monitor error.
% Direct cases run only articulation; service cases run all seven checks.
\begin{tabular}{lccc}
\hline
Case & Original E/A & Repaired E/A & Status \\
\hline
\multicolumn{4}{l}{Direct compiled articulation monitor (18 cases)} \\
\hline
final\_null\_staccato & 2/2 & 0/0 & pass \\
final\_null\_tenuto & 2/2 & 0/0 & pass \\
final\_null\_legato & 2/2 & 0/0 & pass \\
final\_null\_unmarked & 0/0 & 0/0 & pass \\
staccato\_374ms & 0/0 & 0/0 & pass \\
staccato\_375ms & 0/0 & 0/0 & pass \\
staccato\_376ms & 1/1 & 1/1 & pass \\
earlier\_bad\_then\_null\_staccato & 1/1 & 1/1 & pass \\
tenuto\_374ms & 1/1 & 1/1 & pass \\
tenuto\_375ms & 0/0 & 0/0 & pass \\
tenuto\_376ms & 0/0 & 0/0 & pass \\
earlier\_bad\_then\_null\_tenuto & 1/1 & 1/1 & pass \\
legato\_374ms & 1/1 & 1/1 & pass \\
legato\_375ms & 0/0 & 0/0 & pass \\
legato\_376ms & 0/0 & 0/0 & pass \\
earlier\_bad\_then\_null\_legato & 1/1 & 1/1 & pass \\
malformed\_nonnull\_gap & 2/2 & 2/2 & pass \\
malformed\_duration & 2/2 & 2/2 & pass \\
\hline
\multicolumn{4}{l}{Complete converter and service path (11 cases)} \\
\hline
final\_null\_staccato & 2/2 & 0/0 & pass \\
final\_null\_tenuto & 2/2 & 0/0 & pass \\
final\_null\_unmarked & 0/0 & 0/0 & pass \\
earlier\_bad\_then\_null\_staccato & 1/1 & 1/1 & pass \\
staccato\_374ms & 0/0 & 0/0 & pass \\
staccato\_375ms & 0/0 & 0/0 & pass \\
staccato\_376ms & 1/1 & 1/1 & pass \\
earlier\_bad\_then\_null\_tenuto & 1/1 & 1/1 & pass \\
tenuto\_374ms & 1/1 & 1/1 & pass \\
tenuto\_375ms & 0/0 & 0/0 & pass \\
tenuto\_376ms & 0/0 & 0/0 & pass \\
\hline
\end{tabular}
% All six other checks accept for every original/repaired service case.
% Null-gap acceptance exempts this requirement; it does not verify final duration.
% Literal legato has direct-monitor coverage only.
}
\end{center}
\end{document}
